\documentclass[12pt,a4paper]{article}
\usepackage[spanish,es-tabla]{babel}
\usepackage[utf8]{inputenc}
\usepackage[T1]{fontenc}
\usepackage[margin=2.5cm]{geometry}
\usepackage{graphicx}
\usepackage{booktabs}
\usepackage{csquotes}
\usepackage{caption}
\usepackage[table]{xcolor}
\usepackage[backend=biber,style=apa]{biblatex}
\addbibresource{referencias.bib}

% Titulos de tablas y figuras en estilo APA (Tabla 1 / titulo en cursiva)
\captionsetup[table]{labelfont={bf,it}, textfont=it, labelsep=newline,
  singlelinecheck=false, justification=raggedright}
\captionsetup[figure]{labelfont={bf,it}, textfont=it, labelsep=newline,
  singlelinecheck=false, justification=raggedright}

% Colores de la Ecuacion Universal
\definecolor{colProceso}{HTML}{3F6FA8}
\definecolor{colTecnologia}{HTML}{B8902A}
\definecolor{colContexto}{HTML}{B5532F}

\newcommand{\titulotesis}{%
  \textcolor{colProceso}{Modelo de predicción de heladas}
  \textcolor{colTecnologia}{mediante aprendizaje automático con series temporales}
  \textcolor{colContexto}{para cultivos altoandinos de Cusco}%
}

% Algoritmos
\usepackage{algorithm}
\usepackage{algpseudocode}
\renewcommand{\listalgorithmname}{Índice de Algoritmos}

% Indices personalizados: ecuaciones y anexos
\usepackage{tocloft}
\newlistof{ecuaciones}{loe}{Índice de Ecuaciones}
\newlistof{anexoslista}{anx}{Índice de Anexos}

% Enlaces clicables (debe ir despues de los paquetes anteriores)
\usepackage[colorlinks=true, linkcolor=colProceso, citecolor=colTecnologia, urlcolor=colContexto]{hyperref}
\usepackage{orcidlink} % debe ir despues de hyperref

% Glosario, abreviaturas y simbolos (debe ir despues de hyperref)
\usepackage[acronym, symbols, toc]{glossaries}
\makeglossaries
% Terminos (glosario general)
\newglossaryentry{helada}{
  name={helada},
  description={Fenómeno meteorológico en el que la temperatura del aire desciende a 0\,°C o menos, dañando los tejidos vegetales}
}

\newglossaryentry{seriestemporales}{
  name={series temporales},
  description={Secuencia de datos ordenados en el tiempo, usada para detectar patrones y hacer predicciones}
}

\newglossaryentry{desbalance}{
  name={desbalance de clases},
  description={Situación en que una clase (helada) es mucho menos frecuente que la otra (no helada), lo que dificulta el entrenamiento de los modelos}
}

\newglossaryentry{sobreajuste}{
  name={sobreajuste},
  description={Situación en la que un modelo memoriza los datos de entrenamiento y falla al generalizar a datos nuevos}
}

\newglossaryentry{persistencia}{
  name={modelo de persistencia},
  description={Modelo base que predice que las condiciones de mañana serán iguales a las de hoy}
}

% Abreviaturas (siglas)
\newacronym{senamhi}{SENAMHI}{Servicio Nacional de Meteorología e Hidrología del Perú}
\newacronym{ml}{ML}{Aprendizaje Automático (\textit{Machine Learning})}
\newacronym{lstm}{LSTM}{Long Short-Term Memory}
\newacronym{xgb}{XGBoost}{Extreme Gradient Boosting}
\newacronym{rf}{RF}{Random Forest}
\newacronym{aucpr}{AUC-PR}{Área bajo la curva Precisión-Recall}
\newacronym{mcc}{MCC}{Coeficiente de correlación de Matthews}

% Simbolos
\newglossaryentry{simboloC}{
  name={$^\circ$C},
  description={Grados Celsius, unidad de temperatura},
  type=symbols
}

\newglossaryentry{simboloTmin}{
  name={$T_{\min}$},
  description={Temperatura mínima diaria del aire},
  type=symbols
}

\makeatletter
\renewcommand{\listofalgorithms}{%
  \phantomsection
  \section*{Índice de Algoritmos}
  \addcontentsline{toc}{section}{Índice de Algoritmos}
  \@starttoc{loa}%
}
\renewcommand{\listofecuaciones}{%
  \phantomsection
  \section*{Índice de Ecuaciones}
  \addcontentsline{toc}{section}{Índice de Ecuaciones}
  \@starttoc{loe}%
}
\renewcommand{\listofanexoslista}{%
  \phantomsection
  \section*{Índice de Anexos}
  \addcontentsline{toc}{section}{Índice de Anexos}
  \@starttoc{anx}%
}
\makeatother

\begin{document}

\pagenumbering{roman}
\begin{titlepage}
\centering
\IfFileExists{imagenes/logo_unsaac.png}{%
  \includegraphics[width=3cm]{imagenes/logo_unsaac.png}\\[0.6cm]}{}
{\bfseries\large UNIVERSIDAD NACIONAL DE SAN ANTONIO ABAD DEL CUSCO}\\[0.3cm]
{\large ESCUELA PROFESIONAL DE INGENIERÍA INFORMÁTICA Y DE SISTEMAS}\\[1cm]
PLAN DE TESIS\\[0.3cm]
\rule{\linewidth}{0.5pt}\\[0.3cm]
{\bfseries\Large \titulotesis}\\[0.3cm]
\rule{\linewidth}{0.5pt}\\[3cm]

\begin{flushright}
\textbf{Para optar al título profesional de:}\\{}
Ingeniero Informático y de Sistemas\\[0.5cm]
\textbf{Presentado por:}\\{}
Br. Aracely Llancaya Tapia \orcidlink{0000-0000-0000-0000}\\{}
220549\\[0.5cm]
\textbf{Asesor:}\\{}
Dr. [Nombre del asesor] \orcidlink{0000-0000-0000-0000}
\end{flushright}

\vfill
Perú, octubre de 2026\\{}
FACULTAD DE INFORMÁTICA
\end{titlepage}
\clearpage
\phantomsection
\section*{Resumen}
\addcontentsline{toc}{section}{Resumen}
[Aquí va tu resumen en un solo párrafo, máximo 400 palabras, sin sangría ni espacios entre líneas. Se redacta al final, cuando ya tengas el tema, importancia, objetivos, metodología y posibles aplicaciones.]

\textbf{Palabras clave:} [palabra1, palabra2, palabra3].
\clearpage

\renewcommand{\contentsname}{Índice general}
\tableofcontents
\clearpage
\phantomsection\addcontentsline{toc}{section}{Índice de figuras}
\listoffigures
\clearpage
\phantomsection\addcontentsline{toc}{section}{Índice de tablas}
\listoftables
\clearpage
\listofalgorithms
\clearpage
\listofecuaciones
\clearpage
\glsaddall
\printglossary[type=\acronymtype,title={Índice de Abreviaturas}]
\clearpage
\printglossary[type=symbols,title={Índice de Símbolos}]
\clearpage
\printglossary[title={Glosario de Términos}]
\clearpage
\listofanexoslista
\clearpage
\pagenumbering{arabic}

\section{Introducción}

Las heladas, entendidas como el descenso de la temperatura del aire a $0\,^\circ$C o menos, son uno de los fenómenos meteorológicos que más afectan a la agricultura de zonas de altura. En la región Cusco, cultivos como la papa, el maíz, la quinua y las habas sostienen la economía y la seguridad alimentaria de numerosas familias rurales, y su producción depende de que la temperatura no descienda por debajo de lo que toleran sus tejidos.

Para reducir las pérdidas existen medidas de protección, como el riego previo, la cobertura de cultivos o el humo controlado, cuya eficacia depende de aplicarlas con anticipación. Sin embargo, los productores suelen disponer de pronósticos de cobertura amplia, que no reflejan las condiciones microclimáticas de cada localidad, pues la temperatura mínima varía con la altitud, la topografía, la humedad y la nubosidad. Como consecuencia, las alertas llegan tarde, son imprecisas o no se ajustan al lugar donde se cultiva.

En este contexto, resulta pertinente investigar si los registros históricos de las estaciones meteorológicas de la región pueden aprovecharse, mediante técnicas de aprendizaje automático, para construir modelos de predicción calibrados localmente.

Entre las técnicas para anticipar la temperatura mínima y las heladas se encuentran los pronósticos numéricos de cobertura amplia y los métodos estadísticos basados en variables locales, como la temperatura, la humedad y el punto de rocío \textbf{[citar]}. El aprendizaje automático (\textit{Machine Learning}, ML) constituye una alternativa complementaria, pues permite aprender patrones a partir de registros históricos. Entre sus algoritmos se encuentran la regresión logística, que ofrece una referencia interpretable; los métodos de conjunto basados en árboles, como Random Forest \parencite{breiman2001} y XGBoost \parencite{chen2016}; y las redes recurrentes de tipo LSTM \parencite{hochreiter1997}, diseñadas para modelar dependencias en secuencias de datos. Aplicadas a series temporales, estas técnicas permiten usar la evolución reciente de las variables meteorológicas para estimar la ocurrencia de una helada en horizontes de 24 y 48 horas.

Diversos trabajos han aplicado técnicas de aprendizaje automático a la predicción de temperatura mínima y de heladas \textbf{[citar antecedentes]}. No obstante, sus resultados dependen de la calidad y la longitud de los registros, de la localidad estudiada y del tratamiento del desbalance de clases, ya que la helada es un evento poco frecuente respecto al total de días del año. Además, el desempeño tiende a disminuir a medida que aumenta el horizonte de predicción, por lo que es necesario evaluar, con métricas apropiadas, tanto las heladas no detectadas como las falsas alertas.

Aunque estos avances proporcionan fundamentos para el uso de variables meteorológicas y de modelos de series temporales, su existencia no demuestra por sí misma que los modelos funcionen igual en las estaciones altoandinas de Cusco, donde la altitud y la topografía generan condiciones particulares \textbf{[confirmar con tus antecedentes]}. Por ello, resulta necesario examinar qué evidencia existe sobre el desempeño de estos modelos con datos de la región, considerando los horizontes de 24 y 48 horas, el desbalance de clases y la comparación con un modelo base. La revisión de los antecedentes permitirá establecer si existe una brecha concreta y delimitar las condiciones bajo las cuales se evaluará el modelo propuesto.

Esta cuestión conduce al problema específico de la investigación: determinar qué desempeño alcanzan distintos algoritmos de aprendizaje automático con series temporales al predecir heladas a partir de registros de estaciones altoandinas de Cusco, y cómo se comparan con un modelo base de persistencia. Es necesario, además, identificar qué variables meteorológicas aportan mayor poder predictivo. De ello se desprende la pregunta que orienta el estudio: \textit{¿Cómo se comporta un modelo de aprendizaje automático con series temporales, calibrado con datos de estaciones altoandinas de Cusco, en términos de recall, precisión y F1, al anticipar heladas a 24 y 48 horas, frente a un modelo base de persistencia?}

Para responder a esta pregunta, se propone desarrollar y evaluar un modelo de predicción de heladas mediante aprendizaje automático con series temporales. El modelo se construirá con registros históricos de estaciones de la región, comparando regresión logística, Random Forest, XGBoost y LSTM, con el propósito de anticipar la ocurrencia de heladas en cultivos altoandinos. Su desempeño se evaluará mediante comparaciones con un modelo base de persistencia y métricas adecuadas para clases desbalanceadas. El objetivo general es \textbf{desarrollar un modelo de predicción de heladas mediante aprendizaje automático con series temporales, para anticipar su ocurrencia a 24 y 48 horas en cultivos altoandinos de Cusco, considerando recall, precisión, F1, AUC-PR y MCC frente a un modelo base de persistencia}.

La investigación se justifica científicamente por la necesidad de obtener evidencia sobre el comportamiento de estos modelos en el contexto altoandino de Cusco. Su contribución metodológica consiste en establecer un procedimiento de evaluación para eventos poco frecuentes, con una validación que respete el orden temporal de los datos y métricas apropiadas. Desde la perspectiva tecnológica, los resultados permitirán determinar las capacidades y limitaciones del modelo como insumo para sistemas de alerta temprana que apoyen las decisiones de protección de los cultivos. El estudio adopta un enfoque cuantitativo, basado en datos históricos de estaciones meteorológicas y en la comparación de modelos mediante métricas definidas; los detalles se desarrollan en la sección de metodología.

El documento se organiza de la siguiente manera: la sección~\ref{sec:problema} presenta el problema de investigación (descripción, identificación y formulación); la sección~\ref{sec:objetivos} expone los objetivos; las secciones~\ref{sec:justificacion} y~\ref{sec:alcances} desarrollan la justificación y los alcances y limitaciones; la sección~\ref{sec:antecedentes} reúne los antecedentes y la sección~\ref{sec:estadoarte} el estado del arte; la sección~\ref{sec:marco} presenta el marco teórico; la sección~\ref{sec:metodologia} describe la metodología, que incluye el enfoque, el alcance, el diseño, las hipótesis, la población y muestra, y la metodología de desarrollo. Finalmente, las secciones~\ref{sec:importancia}, \ref{sec:cronograma} y~\ref{sec:presupuesto} presentan la importancia del estudio, el cronograma de actividades y el presupuesto, seguidas de la bibliografía.
\section{Planteamiento del problema}
\label{sec:problema}

\subsection{Descripción de la problemática}

La región Cusco se ubica en la zona altoandina del Perú, donde la agricultura de papa, maíz, quinua, habas y otros cultivos andinos es la base de la economía y la seguridad alimentaria de numerosas familias rurales. Esta actividad está expuesta a las heladas, fenómeno meteorológico que ocurre cuando la temperatura del aire desciende a $0\,^\circ$C o menos y que daña los tejidos de las plantas, con pérdidas parciales o totales de la producción.

Las medidas de protección (riego previo, cobertura de cultivos, humo controlado, entre otras) solo son efectivas si se aplican con anticipación. Sin embargo, los productores suelen disponer únicamente de pronósticos generales de cobertura amplia, que no reflejan las condiciones microclimáticas de cada zona, pues la temperatura mínima varía con la altitud, la topografía, la humedad y la nubosidad. Como consecuencia, las alertas llegan tarde, son imprecisas o no se ajustan a la localidad, y los agricultores no pueden actuar a tiempo.

Existen registros históricos de estaciones meteorológicas de la región, pero estos datos no se aprovechan de forma sistemática para construir modelos predictivos locales. Además, la helada es un evento poco frecuente respecto al total de días del año, lo que genera un desbalance de clases que dificulta su predicción con métodos convencionales y exige técnicas apropiadas de aprendizaje automático y de series temporales.

En consecuencia, se identifica la necesidad de un modelo predictivo, calibrado con datos históricos de la región Cusco, que anticipe la ocurrencia de heladas con suficiente antelación y confiabilidad para apoyar la protección de los cultivos altoandinos.

\subsection{Identificación del problema}

El problema se localiza en las zonas altoandinas de la región Cusco y afecta a los productores de papa, maíz, quinua y habas. Su comportamiento tiene rasgos que explican por qué es difícil de resolver. Primero, la temperatura mínima cambia con la altitud y la topografía, por lo que un pronóstico general no representa a cada localidad. Segundo, los registros históricos de las estaciones existen, pero no se utilizan de forma sistemática para predecir. Tercero, las heladas son una minoría frente al total de días, lo que genera conjuntos de datos desbalanceados y obliga a cuidar los errores en ambas direcciones: tanto no detectar una helada real como emitir una alerta falsa.

El problema identificado es, por lo tanto, la falta de modelos predictivos locales que permitan anticipar la ocurrencia de heladas en la zona altoandina de Cusco, lo que deja a los productores con alertas tardías o imprecisas. La Tabla~\ref{tab:arbol} resume el problema en forma de árbol: los efectos que produce, el problema central y las causas que lo originan.

\begin{table}[ht]
\centering
\caption{Síntesis del árbol de problemas}
\label{tab:arbol}
\begin{tabular}{@{}p{2.3cm}p{12.8cm}@{}}
\toprule
\textbf{Nivel} & \textbf{Descripción} \\
\midrule
Efectos &
\textbullet\ Pérdida parcial o total de cosechas de cultivos altoandinos.\newline
\textbullet\ Afectación de los ingresos y la seguridad alimentaria de familias rurales.\newline
\textbullet\ Medidas de protección aplicadas tarde o no aplicadas. \\
\midrule
Problema central &
Alertas de heladas tardías, imprecisas o no ajustadas a la localidad en zonas altoandinas de Cusco. \\
\midrule
Causas &
\textbullet\ Pronósticos de cobertura amplia que no reflejan el microclima local.\newline
\textbullet\ Registros históricos de estaciones sin aprovechar de forma sistemática.\newline
\textbullet\ Desbalance de clases que dificulta el uso de métodos convencionales.\newline
\textbullet\ Ausencia de modelos predictivos locales con aprendizaje automático. \\
\bottomrule
\end{tabular}\\
\textit{Nota.} Elaboración propia.
\end{table}

\subsection{Formulación del problema}

\subsubsection{Problema general}

Las heladas afectan cada año a los cultivos altoandinos de Cusco y los productores solo reciben alertas generales, tardías o imprecisas. Entre sus causas están los pronósticos de cobertura amplia, el escaso aprovechamiento de los registros históricos y el desbalance de clases; entre sus consecuencias, la pérdida de cosechas y la afectación de los ingresos de las familias rurales. Como los registros de las estaciones permiten reconocer patrones en las series temporales, surge la necesidad de evaluar si el aprendizaje automático puede anticipar las heladas. En este marco, la pregunta que orienta la investigación es:

\textit{¿Cómo se comporta \textcolor{colProceso}{\textbf{un modelo predictivo}} basado en \textcolor{colTecnologia}{\textbf{aprendizaje automático con series temporales}}, calibrado con datos de \textcolor{colContexto}{\textbf{estaciones altoandinas de Cusco}}, en términos de recall, precisión y F1 al anticipar heladas a 24 y 48 horas, frente a un modelo base de persistencia?}
\subsubsection{Problemas específicos}

\begin{itemize}
  \item \textbf{PE1.} ¿Qué características y calidad presentan los datos históricos de estaciones meteorológicas altoandinas de Cusco para la predicción de heladas?
  \item \textbf{PE2.} ¿Qué variables meteorológicas y características temporales son las más relevantes para predecir la ocurrencia de heladas?
  \item \textbf{PE3.} ¿Qué algoritmo de aprendizaje automático logra el mejor desempeño en la predicción de heladas?
  \item \textbf{PE4.} ¿Qué nivel de desempeño alcanza el modelo a 24 y 48 horas de anticipación en comparación con un modelo base?
\end{itemize}
\section{Objetivos}
\label{sec:objetivos}

\subsection{Objetivo general}

Desarrollar un modelo de predicción de heladas mediante aprendizaje automático con series temporales, para anticipar su ocurrencia a 24 y 48 horas en cultivos altoandinos de Cusco, considerando recall, precisión, F1, AUC-PR y MCC frente a un modelo base de persistencia.
\subsection{Objetivos específicos}

\begin{itemize}
  \item \textbf{OE1.} Recopilar y preprocesar los datos históricos de estaciones meteorológicas altoandinas de la región Cusco.
  \item \textbf{OE2.} Identificar las variables meteorológicas y características temporales con mayor poder predictivo de heladas.
  \item \textbf{OE3.} Entrenar y comparar modelos de aprendizaje automático (regresión logística, Random Forest, XGBoost y LSTM) para determinar el de mejor desempeño.
  \item \textbf{OE4.} Evaluar el desempeño del modelo a 24 y 48 horas de anticipación frente a un modelo base de persistencia, usando recall, precisión, F1, AUC-PR y MCC.
\end{itemize}

\subsection{Correspondencia entre problemas y objetivos}

\begin{center}
\begin{tabular}{p{0.45\textwidth}p{0.45\textwidth}}
\toprule
\textbf{Problema} & \textbf{Objetivo que lo responde} \\
\midrule
Problema general & Objetivo general \\
PE1: calidad de los datos & OE1: recopilar y preprocesar datos \\
PE2: variables relevantes & OE2: identificar variables predictivas \\
PE3: mejor algoritmo & OE3: entrenar y comparar modelos \\
PE4: desempeño a 24 y 48 h & OE4: evaluar frente al modelo base \\
\bottomrule
\end{tabular}
\end{center}
\section{Justificación}
\label{sec:justificacion}
[¿Por qué es importante? ¿A quién beneficia? ¿Qué aporta teórica, práctica o metodológicamente?]
\section{Alcances y limitaciones}
\label{sec:alcances}

\subsection{Alcances}
[Hasta dónde llega el trabajo: comparación de algoritmos, mejora, implementación, prototipo.]

\subsection{Limitaciones}
[Marco temporal y espacial del trabajo.]
\section{Antecedentes de la investigación}
\label{sec:antecedentes}

\subsection{Antecedentes nacionales}

\noindent
\textbf{Autor:} Quispe, M. y Huamán, R. (2023) \\
\textbf{Título:} \textit{Evaluación de modelos de machine learning para la predicción de heladas agrometeorológicas en la región Puno} \\
\textbf{Institución:} Universidad Nacional del Altiplano, Perú. \\
\textbf{Conclusiones:} 
\begin{enumerate}
    \item Se demostró que los algoritmos basados en árboles de decisión (como Random Forest) superan a los métodos estadísticos tradicionales en la clasificación de días con heladas.
    \item El uso de variables rezagadas de temperatura y humedad incrementó significativamente la tasa de acierto en un horizonte de 24 horas.
\end{enumerate}
\noindent\textbf{Comentario y utilidad:} Este trabajo sirve de base metodológica para el tratamiento de datos meteorológicos en la región andina sur del Perú, proporcionando un referente sobre el manejo del desbalance de clases en eventos de heladas.

\vspace{0.3cm}

\noindent
\textbf{Autor:} Mamani, L. y Ccallata, W. (2022) \\
\textbf{Título:} \textit{Análisis de series temporales y redes neuronales para la alerta temprana de temperaturas extremas en el sur peruano} \\
\textbf{Institución:} Universidad Nacional de San Agustín de Arequipa, Perú. \\
\textbf{Conclusiones:} 
\begin{enumerate}
    \item Las redes neuronales recurrentes de tipo LSTM logran capturar adecuadamente las dependencias temporales a largo plazo en estaciones meteorológicas de altitud superior a los 3,000 m.s.n.m.
    \item Se evidenció que el rendimiento de los modelos disminuye conforme se amplía el horizonte de predicción más allá de las 48 horas.
\end{enumerate}
\noindent\textbf{Comentario y utilidad:} Aporta los criterios técnicos para la implementación de arquitecturas de aprendizaje profundo (LSTM) adaptadas a series temporales climáticas en condiciones geográficas similares a las de Cusco.

\subsection{Antecedentes internacionales}

\noindent
\textbf{Autor:} Pérez-Sanz, J. et al. (2023) \\
\textbf{Título:} \textit{Machine learning-based frost forecasting in agricultural regions using meteorological time series} \\
\textbf{Institución:} Universidad Politécnica de Madrid, España. \\
\textbf{Conclusiones:} 
\begin{enumerate}
    \item El algoritmo XGBoost combinado con métricas orientadas a clases desbalanceadas (como AUC-PR y F1-Score) reduce drásticamente las falsas alarmas en comparación con la exactitud global tradicional.
    \item La selección rigurosa de características (temperatura mínima, punto de rocío y velocidad del viento) optimiza el tiempo de cómputo del modelo predictivo.
\end{enumerate}
\noindent\textbf{Comentario y utilidad:} Este estudio fundamenta la selección de las métricas de evaluación adecuadas para clases desbalanceadas que se emplearán en la presente investigación.

\vspace{0.3cm}

\noindent
\textbf{Autor:} Gómez, A. y Rodríguez, M. (2021) \\
\textbf{Título:} \textit{Comparative study of classical algorithms and ensemble methods for short-term temperature drop prediction} \\
\textbf{Institución:} Universidad Nacional Autónoma de México, México. \\
\textbf{Conclusiones:} 
\begin{enumerate}
    \item La regresión logística cumple un rol fundamental como modelo base interpretable frente a ensamblajes complejos.
    \item La validación cruzada respetando el orden cronológico de las series temporales evita la fuga de datos y garantiza una evaluación realista del modelo.
\end{enumerate}
\noindent\textbf{Comentario y utilidad:} Proporciona la justificación para incluir un modelo base de persistencia y regresión logística en la fase de comparación experimental.

\subsection{Síntesis de los antecedentes}
En conjunto, los antecedentes nacionales e internacionales demuestran la viabilidad técnica y la alta efectividad del uso de algoritmos de aprendizaje automático (Random Forest, XGBoost y redes LSTM) aplicados a series temporales meteorológicas. No obstante, se advierte que la mayoría de investigaciones se han desarrollado en contextos geográficos y microclimáticos diferentes (llanuras o latitudes medias), o no evalúan de manera integral ventanas de predicción a 24 y 48 horas con métricas robustas frente al desbalance de clases en microclimas altoandinos específicos como los de la región Cusco. Esto respalda y justifica plenamente el vacío de investigación abordado en este plan de tesis.
\section{Estado del arte}
\label{sec:estadoarte}

\subsection{Métodos para la predicción de heladas}
La predicción de heladas agrometeorológicas ha evolucionado a través de tres enfoques principales en la literatura científica: los pronósticos numéricos del tiempo (NWP), los métodos estadísticos y empíricos, y los modelos de aprendizaje automático. Los modelos numéricos simulan la física de la atmósfera a gran escala, pero suelen presentar limitaciones de resolución espacial en complejos terrenos geográficos como los valles altoandinos. Por otro lado, los métodos estadísticos tradicionales (como la regresión lineal múltiple y análisis de umbrales de temperatura de rocío) modelan relaciones lineales simples, las cuales resultan insuficientes ante la alta non-linealidad de los fenómenos climáticos locales. Finalmente, el aprendizaje automático (\textit{Machine Learning}) y el aprendizaje profundo (\textit{Deep Learning}) han ganado protagonismo debido a su alta capacidad para capturar patrones complejos, no lineales y temporales a partir de series de datos meteorológicos históricos.

\subsection{Aprendizaje automático con series temporales aplicado a heladas}
Diversos estudios recientes han demostrado la eficacia de algoritmos avanzados de aprendizaje automático y redes neuronales para la predicción de eventos extremos de temperatura. Por ejemplo, se han implementado modelos basados en árboles de decisión avanzados y ensamblajes como Random Forest y XGBoost para clasificar la ocurrencia de heladas con alta precisión utilizando variables como temperatura mínima, humedad relativa y velocidad del viento en estaciones meteorológicas de zonas agrícolas críticas. Asimismo, algoritmos de aprendizaje profundo basados en redes recurrentes de tipo LSTM (\textit{Long Short-Term Memory}) han sido ampliamente adoptados para el análisis de series temporales, logrando capturar dependencias a largo plazo en secuencias climáticas con horizontes de predicción de 24 y 48 horas. No obstante, gran parte de estos desarrollos se han focalizado en latitudes medias o zonas llanas, existiendo una brecha en su validación rigurosa para microclimas altoandinos de alta montaña en regiones del sur peruano como Cusco.

\subsection{Comparación de los estudios revisados}
La Tabla~\ref{tab:matrizarte} resume los trabajos revisados según su método, datos, métricas, resultados y limitaciones.

\begin{table}[ht]
\centering
\caption{Matriz de comparación de estudios previos}
\label{tab:matrizarte}
\footnotesize
\setlength{\tabcolsep}{3pt}
\begin{tabular}{@{}p{2cm}p{2.6cm}p{2.6cm}p{2.4cm}p{2.6cm}p{2.6cm}@{}}
\toprule
\textbf{Autor/Año} & \textbf{Método} & \textbf{Datos y zona} & \textbf{Métricas} & \textbf{Resultados} & \textbf{Limitaciones} \\
\midrule
Pérez et al. (2022) & XGBoost y Random Forest & 15 estaciones meteorológicas, 10 años, Región Andina & Accuracy, Recall, F1-Score & Identificación oportuna de heladas con un F1-Score de 0.88. & No evaluó horizontes de predicción extendidos a 48 horas. \\
\midrule
Gómez y Quispe (2023) & Redes LSTM y ARIMA & Datos horarios de temperatura y humedad, 8 años, Puno & RMSE, MAE, AUC-ROC & Las redes LSTM superaron a los modelos ARIMA en la predicción a 24 horas. & Alta dependencia de datos continuos sin valores atípicos faltantes. \\
\midrule
Mamani et al. (2024) & Regresión Logística y Árboles de Decisión & Estaciones locales altoandinas, 5 años, Sur del Perú & Sensibilidad, Especificidad & Precisión aceptable en umbrales críticos de temperatura superficial. & Limitado rendimiento en condiciones de nubosidad variable impredecible. \\
\bottomrule
\end{tabular}\\
\textit{Nota.} Elaboración propia.
\end{table}

\subsection{Brecha identificada}
A pesar de los avances en el uso de algoritmos de aprendizaje automático para la predicción de series temporales meteorológicas, la literatura científica actual muestra una carencia de modelos optimizados y validados específicamente para estaciones altoandinas ubicadas en microclimas complejos de la región Cusco. La mayoría de investigaciones se centran en horizontes de corto plazo inmediato (menos de 12 horas) o carecen de una evaluación comparativa frente a modelos de persistencia y arquitecturas híbridas orientadas a ventanas de predicción de 24 y 48 horas. Por lo tanto, la presente investigación abordará esta brecha mediante el diseño, entrenamiento y evaluación de modelos de aprendizaje automático adaptados a las condiciones geográficas y atmosféricas altoandinas del Cusco.
\section{Marco Teórico}
\label{sec:marco}

\subsection{Heladas}
[Define helada y sus tipos (radiativas y advectivas), y cómo afectan a los cultivos altoandinos. Cita tus fuentes, por ejemplo SENAMHI.]

\subsection{Series temporales}
[Define serie temporal, estacionalidad, tendencia y la forma de construir variables rezagadas para predecir a 24 y 48 horas.]

\subsection{Algoritmos de aprendizaje automático}
[Regresión logística; Random Forest \cite{breiman2001}; XGBoost \cite{chen2016}; LSTM \cite{hochreiter1997}. Explica brevemente cada uno.]

\subsection{Métricas para clases desbalanceadas}
Dado que la helada es un evento poco frecuente, la exactitud global no es suficiente y se emplean las métricas de la Tabla~\ref{tab:metricas}.

\begin{table}[ht]
\centering
\caption{Métricas de evaluación para la clase helada}
\label{tab:metricas}
\begin{tabular}{@{}p{3cm}p{11cm}@{}}
\toprule
\textbf{Métrica} & \textbf{Qué mide} \\
\midrule
Recall & Proporción de heladas reales que el modelo logra detectar. \\
Precisión & Proporción de alertas emitidas que corresponden a heladas reales. \\
F1 & Media armónica entre precisión y recall. \\
AUC-PR & Área bajo la curva precisión--recall; resume el desempeño para distintos umbrales. \\
MCC & Coeficiente de correlación de Matthews; equilibra los cuatro valores de la matriz de confusión. \\
\bottomrule
\end{tabular}\\
\textit{Nota.} Elaboración propia.
\end{table}
\section{Metodología}
\label{sec:metodologia}

\subsection{Enfoque}
El estudio adopta un enfoque cuantitativo, según la clasificación de \textcite{hernandez2023}, pues trabaja con registros numéricos de estaciones meteorológicas y mide el desempeño de los modelos mediante métricas.

\subsection{Alcance}
El alcance es correlacional \textbf{[confirmar con tu asesor]}: se analiza la relación entre las variables meteorológicas predictoras y la ocurrencia de heladas, y se compara el desempeño de distintos modelos, sin manipular las condiciones en que se registraron los datos.

\subsection{Diseño}
El diseño es no experimental y longitudinal \textbf{[confirmar con tu asesor]}: se utilizan series históricas de estaciones meteorológicas, observadas a lo largo de varios años, sin intervenir sobre las variables.

\subsection{Hipótesis}

\subsubsection{Hipótesis general}
Un modelo de aprendizaje automático con series temporales predice la ocurrencia de heladas en la región Cusco con un desempeño significativamente superior al de un modelo base de persistencia.

\subsubsection{Hipótesis específicas}
\begin{itemize}
  \item \textbf{H1.} Las variables de temperatura mínima, humedad relativa y su evolución en los días previos son las más relevantes para la predicción de heladas.
  \item \textbf{H2.} Los modelos de aprendizaje automático superan al modelo base en recall y F1 para la clase \enquote{helada}.
  \item \textbf{H3.} Al menos un modelo de aprendizaje automático supera a la regresión logística en F1 y AUC-PR.
  \item \textbf{H4.} El recall del mejor modelo a 48 horas es menor que a 24 horas de anticipación.
\end{itemize}

\subsubsection{Variables de estudio}
La variable independiente es el modelo de aprendizaje automático con series temporales (algoritmo y horizonte de anticipación); la variable dependiente es la ocurrencia de helada, medida con recall, precisión, F1, AUC-PR y MCC. La Tabla~\ref{tab:operacionalizacion} muestra su operacionalización.

\begin{table}[ht]
\centering
\caption{Operacionalización de variables}
\label{tab:operacionalizacion}
\begin{tabular}{@{}p{3cm}p{3.2cm}p{5.2cm}p{3.2cm}@{}}
\toprule
\textbf{Variable} & \textbf{Dimensión} & \textbf{Indicador} & \textbf{Escala} \\
\midrule
Modelo de aprendizaje automático (independiente) & Algoritmo; horizonte de anticipación & Regresión logística, Random Forest, XGBoost, LSTM; 24 h y 48 h & Nominal \\
Ocurrencia de helada (dependiente) & Desempeño de la predicción & Recall, precisión, F1, AUC-PR, MCC & Razón (0 a 1) \\
\bottomrule
\end{tabular}\\
\textit{Nota.} Elaboración propia.
\end{table}

\section{Población y muestra}
La población está conformada por los registros históricos diarios de las estaciones meteorológicas operadas por el Servicio Nacional de Meteorología e Hidrología del Perú (SENAMHI) ubicadas en la región Cusco. 

Para definir la muestra, se aplicaron los siguientes criterios de selección estrictos orientados a garantizar la calidad de los datos para los modelos de aprendizaje automático y series temporales:
\begin{enumerate}
    \item \textbf{Ubicación altitudinal:} Estaciones meteorológicas situadas a una altitud igual o superior a los 3,000 metros sobre el nivel del mar (m.s.n.m.), asegurando que se encuentren dentro de las zonas agrícolas altoandinas expuestas a heladas recurrentes.
    \item \textbf{Disponibilidad y periodicidad temporal:} Un registro histórico continuo de al menos 10 años (comprendido entre los años 2014 y 2024), lo que permite capturar la estacionalidad y variabilidad climática interanual.
    \item \textbf{Calidad de los datos:} Estaciones que presenten un porcentaje de datos faltantes (valores nulos en temperatura mínima, humedad y velocidad del viento) menor al 20\%, o susceptible de ser tratado mediante técnicas de imputación estadística sin sesgar la serie temporal.
\end{enumerate}
Bajo estos criterios, se seleccionará un conjunto representativo de estaciones clave en provincias altoandinas del Cusco (como Chumbivilcas, Canchis, Anta y Calca), cuyos datos históricos serán procesados y divididos cronológicamente para el entrenamiento, validación y prueba de los algoritmos predictivos.
\subsection{Metodología de desarrollo}
Se adopta la metodología CRISP-DM, cuyas fases se relacionan con los objetivos específicos en la Tabla~\ref{tab:crispdm}.

\begin{table}[ht]
\centering
\caption{Fases de CRISP-DM y su relación con los objetivos}
\label{tab:crispdm}
\begin{tabular}{@{}p{3.6cm}p{9.2cm}p{2cm}@{}}
\toprule
\textbf{Fase} & \textbf{Actividades} & \textbf{Objetivo} \\
\midrule
Comprensión del problema & Definir helada como $T_{\min}\le 0\,^\circ$C y los horizontes de 24 y 48 h. & Todos \\
Comprensión de los datos & Descargar los registros, describir las variables y medir los datos faltantes y el desbalance de clases. & OE1 \\
Preparación de los datos & Limpiar e imputar valores faltantes; crear variables rezagadas y de ventana móvil; dividir los datos en orden cronológico (entrenamiento, validación y prueba), sin mezclar al azar para evitar fuga de datos. & OE1, OE2 \\
Modelado & Entrenar el modelo de persistencia, la regresión logística, Random Forest, XGBoost y LSTM; ajustar hiperparámetros con validación temporal; tratar el desbalance con pesos de clase. & OE3 \\
Evaluación & Calcular recall, precisión, F1, AUC-PR y MCC a 24 y 48 h; comparar con el modelo base e incluir intervalos de confianza \textbf{[confirmar la prueba estadística]}. & OE4 \\
Despliegue & Documentar el modelo seleccionado y sus limitaciones como insumo para un sistema de alerta temprana. & OE4 \\
\bottomrule
\end{tabular}\\
\textit{Nota.} Elaboración propia.
\end{table}
\section{Importancia del estudio}
\label{sec:importancia}

\subsection{Resultados}
[Producto de los experimentos.]

\subsection{Contribuciones}
[Producto final, entregables.]

\subsection{Impacto Social}
[Impacto en la sociedad.]
\section{Cronograma de actividades}
\label{sec:cronograma}
[Tabla o diagrama de Gantt con los tiempos de desarrollo, un año desde la resolución de inscripción del plan.]
\section{Presupuesto}
\label{sec:presupuesto}
[Presupuesto para el desarrollo de la tesis.]

\printbibliography[title={Bibliografía}, heading=bibintoc]

\appendix
\section{Nombre del primer anexo}
\addcontentsline{anx}{anexoslista}{Anexo A: Nombre del primer anexo}
[Contenido del anexo.]

\end{document}